% !TEX root = ../../main.tex

\subsection{From Test Annotations to Final Clusters}

At this point both inferential questions have already been answered, so the
remaining task is no longer statistical calibration but decision extraction.

Once the edge-test and sibling-test annotations are available, the inferential
work is finished. The remaining step is purely algorithmic: traverse the tree
from the root downward and decide where the splits are kept. The traversal does
not estimate any new model or compute any new $p$-values; it turns the
precomputed evidence tables into a final partition.

The final clusters are produced by a top-down traversal that reads the two test
annotation tables. A tree break means that TBS accepts the proposed split at an
internal node and descends into its two children. For an internal node $u$ with
children $l(u)$ and $r(u)$, the break rule is
\begin{equation}
B_{\mathrm{TBS}}(u) = 1
\end{equation}
if and only if all three conditions hold:
\begin{enumerate}[leftmargin=1.5em]
  \item $u$ has exactly two children;
  \item $S^{\mathrm{edge}}(l(u)) = 1$ or $S^{\mathrm{edge}}(r(u)) = 1$;
  \item $S^{\mathrm{sib}}(u) = 1$.
\end{enumerate}
If these conditions hold, traversal continues into both children. Otherwise the
method would ordinarily stop and declare $u$ a terminal cluster.

The implementation documented here uses a pass-through traversal. Let
$D^{\mathrm{sib}}(u)$ be the indicator that some proper descendant of $u$ has a
significant sibling split among nodes for which the sibling test was run. Then
the traversal still descends into
both children when
\begin{equation}
B_{\mathrm{TBS}}(u) = 0,
\qquad
u \text{ has exactly two children},
\qquad
S^{\mathrm{edge}}(l(u)) \vee S^{\mathrm{edge}}(r(u)) = 1,
\qquad
D^{\mathrm{sib}}(u) = 1.
\end{equation}
This pass-through rule is used to prevent a shallow sibling-test failure from
masking deeper structure. The terminal nodes returned by this traversal define
\(\mathcal C_{\mathrm{TBS}}\), the final cluster partition.
